\section{Appendix}\label{sec:Appendix}

\subsection{Base Localization Node Implementation}\label{subsec:Base implementation}
\begin{verbatim}
#!/usr/bin/env python3

from zlib import crc32
from numpy import ndarray, array, linspace, isinf, isnan, diag, deg2rad, zeros
from numpy.random import default_rng

from rclpy.node import Node
from tf2_ros import Buffer, TransformBroadcaster, TransformListener

from nav_msgs.msg import Odometry, OccupancyGrid
from geometry_msgs.msg import PoseWithCovarianceStamped, PointStamped
from sensor_msgs.msg import LaserScan

from devol_localization.utils import euler_to_quaternion


__author__ = 'Jacob Taylor Cassady'
__email__ = 'jcassad1@jh.edu'


class BaseLocalizationNode(Node):
    def __init__(
        self,
        node_name,
        pose_topic: str = '/devol_drive/pose',
        runtime_topic: str = '/devol_drive/pose/runtime',
    ):
        super().__init__(node_name)

        # Declare Parameters
        self.declare_parameter('odom_topic', '/devol_drive/odom')
        self.declare_parameter('map_topic', '/devol_drive/localization/projected_map')
        self.declare_parameter('scan_topic', '/devol_drive/sensors/lidar2d_0/scan')
        self.declare_parameter('pose_topic', pose_topic)
        self.declare_parameter('runtime_topic', runtime_topic)
        self.declare_parameter('map_frame', 'map')
        self.declare_parameter('odom_frame', 'odom')
        self.declare_parameter('base_frame', 'a200_base_link')
        self.declare_parameter('seed', '31')

        # Retrieve Parameters
        self.odom_topic: str = self.get_parameter('odom_topic').value
        self.map_topic: str = self.get_parameter('map_topic').value
        self.scan_topic: str = self.get_parameter('scan_topic').value
        self.pose_topic: str = self.get_parameter('pose_topic').value
        self.runtime_topic: str = self.get_parameter('runtime_topic').value
        self.map_frame: str = self.get_parameter('map_frame').value
        self.odom_frame: str = self.get_parameter('odom_frame').value
        self.base_frame: str = self.get_parameter('base_frame').value
        self.seed: int = int(self.get_parameter('seed').value)

        # Shared State and Models
        self.v: float = 0.0
        self.omega: float = 0.0
        self.x_est: ndarray = array([0.0, 0.0, 0.0])
        self.covariance: ndarray = diag([0.1**2, 0.1**2, deg2rad(5.0) ** 2])

        # Proportional Control Noise Parameters
        self.alpha_tv_causedby_tv: float = 0.10
        self.alpha_tv_causedby_rv: float = 0.01
        self.alpha_rv_causedby_tv: float = 0.01
        self.alpha_rv_causedby_rv: float = 0.20

        # Initialize random number generator
        self.rng = default_rng(self.seed)

        self.latest_odom_pose = array([0.0, 0.0, 0.0])
        self.last_map_crc = None
        self.last_scan_time = None

        # Shared Publishers
        self.pub_pose = self.create_publisher(PoseWithCovarianceStamped, self.pose_topic, 10)
        self.runtime_pub = self.create_publisher(PointStamped, self.runtime_topic, 10)

    def odom_callback(self, msg: Odometry):
        """Base callback to track the latest odometry twist."""
        self.v: float = msg.twist.twist.linear.x
        self.omega: float = msg.twist.twist.angular.z
        self.u: ndarray = array([self.v, self.omega])

    def is_map_new(self, msg: OccupancyGrid) -> bool:
        """Computes a CRC32 hash of the map data.
        Returns True if the map is new/changed, False if it is identical."""
        # Convert the map data array to continuous bytes and hash it
        current_crc = crc32(bytes(msg.data))

        if self.last_map_crc == current_crc:
            return False

        self.last_map_crc = current_crc
        return True

    def get_dt_from_header(self, header):
        """Calculates the time delta (dt) from a message header.
        Updates self.last_scan_time.
        Returns dt (float) if valid, or None if the loop should be skipped."""
        msg_time = header.stamp.sec + (header.stamp.nanosec / 1e9)

        if self.last_scan_time is None:
            self.last_scan_time = msg_time
            return None

        dt = msg_time - self.last_scan_time

        # Prevent massive dt jumps if processing fell way behind
        if dt > 1.0:
            self.get_logger().warn(f'Large dt detected ({dt:.2f}s). Resetting.')
            self.last_scan_time = msg_time
            return None

        self.last_scan_time = msg_time
        return dt

    @staticmethod
    def preprocess_scan(msg: LaserScan):
        """Standardizes LiDAR extraction for all algorithms."""
        angles: ndarray = linspace(msg.angle_min, msg.angle_max, len(msg.ranges))
        ranges: ndarray = array(msg.ranges)

        # Filter out infinite/bad rays
        eps = 0.02
        valid = (
            (ranges > msg.range_min)
            & (ranges < (msg.range_max - eps))
            & ~isinf(ranges)
            & ~isnan(ranges)
        )
        return ranges[valid], angles[valid]

    def get_laser_offset(self):
        """Looks up the physical LiDAR mount offset."""
        return array([0.32825, 0.0, 0.0])

    @staticmethod
    def control_noise_covariance(
        u: ndarray,
        alpha_tv_causedby_tv: float,
        alpha_tv_causedby_rv: float,
        alpha_rv_causedby_tv: float,
        alpha_rv_causedby_rv: float,
    ) -> ndarray:
        """M"""
        v: float = u[0]
        omega: float = u[1]

        sigma_v: float = alpha_tv_causedby_tv * abs(v) + alpha_tv_causedby_rv * abs(omega) + 0.001
        sigma_w: float = alpha_rv_causedby_tv * abs(v) + alpha_rv_causedby_rv * abs(omega) + 0.001

        return diag([sigma_v**2, sigma_w**2])

    @staticmethod
    def update_covariance(A: ndarray, cov: ndarray, R: ndarray) -> ndarray:
        """Sigma"""
        return A @ cov @ A.T + R

    def publish_pose_and_tf_runtime(self, stamp, cov_3x3, start_time: float, end_time: float):
        """Standardized publishing of Pose + Covariance and REP-105 TF."""
        # Publish Pose
        pose_msg = PoseWithCovarianceStamped()
        pose_msg.header.stamp = stamp
        pose_msg.header.frame_id = self.map_frame

        pose_msg.pose.pose.position.x = self.x_est[0]
        pose_msg.pose.pose.position.y = self.x_est[1]
        pose_msg.pose.pose.orientation = euler_to_quaternion(self.x_est[2])

        # Pack 3x3 covariance into ROS 6x6 format
        cov_6x6 = zeros((6, 6))
        cov_6x6[0, 0], cov_6x6[0, 1] = cov_3x3[0, 0], cov_3x3[0, 1]
        cov_6x6[1, 0], cov_6x6[1, 1] = cov_3x3[1, 0], cov_3x3[1, 1]
        cov_6x6[0, 5], cov_6x6[5, 0] = cov_3x3[0, 2], cov_3x3[2, 0]
        cov_6x6[1, 5], cov_6x6[5, 1] = cov_3x3[1, 2], cov_3x3[2, 1]
        cov_6x6[5, 5] = cov_3x3[2, 2]

        pose_msg.pose.covariance = tuple(cov_6x6.flatten())
        self.pub_pose.publish(pose_msg)

        # Publish the runtime as a PointStamped message
        runtime_msg = PointStamped()
        runtime_msg.header.stamp = stamp
        runtime_msg.point.x = end_time - start_time
        self.runtime_pub.publish(runtime_msg)
\end{verbatim}

\newpage
\subsection{Dead Reckoning Node Implementation}\label{subsec:DR implementation}
\begin{verbatim}
#!/usr/bin/env python3

from typing import Tuple, Optional
from time import perf_counter
from numpy import diag, array, ndarray, sin, cos, zeros, deg2rad

import rclpy
from nav_msgs.msg import Odometry

from devol_localization.base_localization_node import BaseLocalizationNode
from devol_localization.motion_models import UnicycleMotionModel
from devol_localization.utils import wrap_angle


__author__ = 'Jacob Taylor Cassady'
__email__ = 'jcassad1@jh.edu'


class DeadReckoningNode(BaseLocalizationNode):
    def __init__(self):
        super().__init__(
            'dead_reckoning_node',
            pose_topic='/devol_drive/dr_pose',
            runtime_topic='/devol_drive/dr_pose/runtime',
        )

        # State variables.
        self.x_est: ndarray = array([0.0, 0.0, 0.0])  # [x, y, theta]
        self.covariance: ndarray = diag([0.1**2, 0.1**2, deg2rad(5.0) ** 2])
        self.localization_runtime = 0.0

        self.create_subscription(Odometry, self.odom_topic, self.odom_callback, 10)

        self.get_logger().info('Dead Reckoning Node Started!')

    def odom_callback(self, msg: Odometry):
        """Overrides base class implementation"""
        start_time: float = perf_counter()
        dt: Optional[float] = self.get_dt_from_header(msg.header)
        if dt is None:
            self.localization_runtime += perf_counter() - start_time
            return

        self.x_est, self.covariance = DeadReckoningNode._predict(
            prev_x=self.x_est,
            prev_cov=self.covariance,
            v=msg.twist.twist.linear.x,
            omega=msg.twist.twist.angular.z,
            dt=dt,
            alpha_tv_causedby_tv=self.alpha_tv_causedby_tv,
            alpha_tv_causedby_rv=self.alpha_tv_causedby_rv,
            alpha_rv_causedby_tv=self.alpha_rv_causedby_tv,
            alpha_rv_causedby_rv=self.alpha_rv_causedby_rv,
        )
        self.localization_runtime += perf_counter() - start_time

        # Publish Pose, TF, and Runtime using the base class
        self.publish_pose_and_tf_runtime(
            stamp=msg.header.stamp,
            cov_3x3=self.covariance,
            start_time=0.0,
            end_time=self.localization_runtime,
        )

        # Reset runtime for the next cycle
        self.localization_runtime = 0.0

    @staticmethod
    def mapped_process_noise(x: ndarray, M: ndarray, dt: float) -> ndarray:
        """R"""
        G: ndarray = UnicycleMotionModel.control_jacobian(x=x, dt=dt)
        return G @ M @ G.T * dt

    @staticmethod
    def _predict(
        prev_x: ndarray,
        prev_cov: ndarray,
        v: float,
        omega: float,
        dt: float,
        alpha_tv_causedby_tv: float,
        alpha_tv_causedby_rv: float,
        alpha_rv_causedby_tv: float,
        alpha_rv_causedby_rv: float,
    ) -> Tuple[ndarray, ndarray]:
        """Prediction step. Apply Kinematic model"""
        u: ndarray = array([v, omega])

        # Apply model forward kinematics
        x: ndarray = UnicycleMotionModel.forward(x=prev_x, u=u, noise=zeros(3), dt=dt)
        x: ndarray = wrap_angle(x)

        # Proportional Control Noise Covariance (M)
        M: ndarray = DeadReckoningNode.control_noise_covariance(
            u=u,
            alpha_tv_causedby_tv=alpha_tv_causedby_tv,
            alpha_tv_causedby_rv=alpha_tv_causedby_rv,
            alpha_rv_causedby_tv=alpha_rv_causedby_tv,
            alpha_rv_causedby_rv=alpha_rv_causedby_rv,
        )

        # Mapped Process Noise (R)
        R: ndarray = DeadReckoningNode.mapped_process_noise(x=x, M=M, dt=dt)

        # Calculate the Jacobian using the previous step.
        A: ndarray = UnicycleMotionModel.state_jacobian(x=prev_x, u=u, dt=dt)

        # Update the covariance
        cov: ndarray = DeadReckoningNode.update_covariance(A=A, cov=prev_cov, R=R)

        return x, cov


def main(args=None):
    rclpy.init(args=args)
    rclpy.spin(DeadReckoningNode())
    rclpy.shutdown()


if __name__ == '__main__':
    main()
\end{verbatim}

\newpage
\subsection{Extended Kalman Filter Node Implementation}\label{subsec:EKF implementation}
\begin{verbatim}
#!/usr/bin/env python3

from typing import Tuple
from time import perf_counter
from scipy.spatial import KDTree
from scipy.stats import chi2
from numpy import (
    ndarray,
    array,
    zeros,
    diag,
    deg2rad,
    mean,
    sum,
    sin,
    cos,
    vstack,
    where,
    int8 as np_int8,
    eye,
    arctan2,
    trace,
    sqrt,
    clip,
)
from numpy.linalg import svd, det, inv

import rclpy
from rclpy.qos import DurabilityPolicy, QoSProfile, ReliabilityPolicy

# ROS 2 Messages
from nav_msgs.msg import OccupancyGrid, Odometry
from sensor_msgs.msg import LaserScan

# Import devol localization classes
from devol_localization.base_localization_node import BaseLocalizationNode
from devol_localization.motion_models import UnicycleMotionModel
from devol_localization.utils import wrap_angle


__author__ = 'Jacob Taylor Cassady'
__email__ = 'jcassad1@jh.edu'


# Chi-square threshold for 3 degrees of freedom at 95% confidence
CHI2_TRESHOLD_3DOF_95CONF: float = 7.815


class EKFLocalizationNode(BaseLocalizationNode):
    def __init__(self):
        super().__init__(
            'ekf_localization_node',
            '/devol_drive/ekf_pose',
            runtime_topic='/devol_drive/ekf_pose/runtime',
        )

        self.map_data = None
        self.W_pose: ndarray = diag([0.05**2, 0.05**2, deg2rad(0.2) ** 2])

        # State variables.
        self.x_est: ndarray = array(
            [0.0, 0.0, 0.0, 1.0]
        )  # [x, y, theta, s], where true yaw rate = s * raw odometry yaw
        self.covariance: ndarray = diag([0.1**2, 0.1**2, deg2rad(5.0) ** 2, 0.3**2])
        self.q_yaw_scale: float = 0.02

        # Time synchronization variables
        self.odom_buffer: list = []
        self.last_time: float = None

        # Total runtime accumulator
        self.localization_runtime: float = 0.0

        # Subscribers / Publishers
        map_qos: QoSProfile = QoSProfile(
            depth=1,
            reliability=ReliabilityPolicy.RELIABLE,
            durability=DurabilityPolicy.TRANSIENT_LOCAL,
        )
        self.create_subscription(OccupancyGrid, self.map_topic, self.map_callback, map_qos)
        self.create_subscription(LaserScan, self.scan_topic, self.scan_callback, 10)
        self.create_subscription(Odometry, self.odom_topic, self.odom_callback, 10)

        self.get_logger().info('EKF Localization Node Started!')

    def odom_callback(self, msg: Odometry):
        """Base callback to track the latest odometry twist."""
        start_time: float = perf_counter()
        # Convert ROS timestamp to a float in seconds
        stamp = msg.header.stamp
        t_odom = stamp.sec + stamp.nanosec * 1e-9
        v = msg.twist.twist.linear.x
        omega = msg.twist.twist.angular.z

        # Add to the buffer
        self.odom_buffer.append((t_odom, v, omega))
        self.localization_runtime += perf_counter() - start_time

    def map_callback(self, msg: OccupancyGrid):
        """Extracts obstacle coordinates and builds KDTree."""
        start_time: float = perf_counter()

        if not self.is_map_new(msg):
            self.localization_runtime += perf_counter() - start_time
            return

        grid: ndarray = array(msg.data, dtype=np_int8).reshape((msg.info.height, msg.info.width))

        # Find the row/col indices of all cells that are likely obstacles
        obstacle_y, obstacle_x = where(grid > 50)

        # Convert indices to real-world metric coordinates
        map_pts_x: ndarray = msg.info.origin.position.x + (obstacle_x + 0.5) * msg.info.resolution
        map_pts_y: ndarray = msg.info.origin.position.y + (obstacle_y + 0.5) * msg.info.resolution

        # Stack into an Nx2 array of obstacle points
        map_points: ndarray = vstack((map_pts_x, map_pts_y)).T

        # Build the KDTree for the ICP algorithm
        self.map_data: dict = {'points': map_points, 'kdtree': KDTree(map_points)}
        self.get_logger().info(
            f'Map received. Built KDTree with {len(map_points)} obstacle points.'
        )

        self.localization_runtime += perf_counter() - start_time

    def scan_callback(self, msg: LaserScan):
        """The main loop of the EKF."""
        start_time: float = perf_counter()

        if self.map_data is None:
            self.localization_runtime += perf_counter() - start_time
            return  # Wait for map to load

        # Get exact timestamp of the laser scan
        stamp = msg.header.stamp
        t_scan = stamp.sec + stamp.nanosec * 1e-9

        # Initialize time tracker on the very first scan
        if self.last_time is None:
            self.last_time = t_scan
            self.localization_runtime += perf_counter() - start_time
            return

        # Process all buffered odom messages stamped AT or BEFORE the scan time
        x_pred, cov_pred = self.x_est.copy(), self.covariance.copy()
        current_v = 0.0
        current_omega = 0.0
        while self.odom_buffer and self.odom_buffer[0][0] <= t_scan:
            t_odom, v, omega = self.odom_buffer.pop(0)

            dt = t_odom - self.last_time
            if dt > 0:
                x_pred, cov_pred = EKFLocalizationNode._predict(
                    prev_x=x_pred,
                    prev_cov=cov_pred,
                    v=v,
                    omega=omega,
                    dt=dt,
                    q_yaw_scale=self.q_yaw_scale,
                    alpha_tv_causedby_tv=self.alpha_tv_causedby_tv,
                    alpha_tv_causedby_rv=self.alpha_tv_causedby_rv,
                    alpha_rv_causedby_tv=self.alpha_rv_causedby_tv,
                    alpha_rv_causedby_rv=self.alpha_rv_causedby_rv,
                )
                self.last_time = t_odom

            # Keep track of the latest velocities for the partial step
            current_v = v
            current_omega = omega

        # Add one final partial step exactly up to the scan timestamp
        dt_partial = t_scan - self.last_time
        if dt_partial > 0:
            x_pred, cov_pred = EKFLocalizationNode._predict(
                prev_x=x_pred,
                prev_cov=cov_pred,
                v=current_v,
                omega=current_omega,
                dt=dt_partial,
                q_yaw_scale=self.q_yaw_scale,
                alpha_tv_causedby_tv=self.alpha_tv_causedby_tv,
                alpha_tv_causedby_rv=self.alpha_tv_causedby_rv,
                alpha_rv_causedby_tv=self.alpha_rv_causedby_tv,
                alpha_rv_causedby_rv=self.alpha_rv_causedby_rv,
            )
            self.last_time = t_scan

        # update step (Using LaserScan ICP)
        z, is_valid, match_error = EKFLocalizationNode.scan_match(
            msg=msg,
            x_pred=x_pred,
            cov_pred=cov_pred,
            map_data=self.map_data,
            laser_offset=self.get_laser_offset(),
        )

        if is_valid:
            self.x_est, self.covariance, accepted = EKFLocalizationNode._update(
                x_pred=x_pred, cov_pred=cov_pred, W_pose=self.W_pose, z=z, match_error=match_error
            )

            if not accepted:
                self.get_logger().debug('ICP match rejected by Mahalanobis check.')
        else:
            self.x_est = x_pred
            self.covariance = cov_pred

        # Publish state
        self.localization_runtime += perf_counter() - start_time
        self.publish_pose_and_tf_runtime(
            stamp=msg.header.stamp,
            cov_3x3=self.covariance,
            start_time=0.0,
            end_time=self.localization_runtime,
        )

        # Reset runtime for the next cycle
        self.localization_runtime = 0.0

    @staticmethod
    def scan_match(
        msg: LaserScan,
        x_pred: ndarray,
        cov_pred: ndarray,
        map_data: dict,
        laser_offset: ndarray,
        max_iter: int = 200,
        tolerance: float = 1e-4,
    ):
        """2D Iterative Closest Point using SVD."""
        # Filter out invalid rays
        ranges, angles = EKFLocalizationNode.preprocess_scan(msg=msg)

        if len(ranges) < 10:
            return array([0.0, 0.0, 0.0, 1.0]), False, 0.0  # Not enough valid points

        # Nx2 array of local laser points
        laser_pts: ndarray = vstack((ranges * cos(angles), ranges * sin(angles))).T

        R_offset: ndarray = array(
            [
                [cos(laser_offset[2]), -sin(laser_offset[2])],
                [sin(laser_offset[2]), cos(laser_offset[2])],
            ]
        )
        local_pts: ndarray = (R_offset @ laser_pts.T).T + laser_offset[:2]

        z: ndarray = x_pred[:3].copy()
        prev_error: float = float('inf')

        # Calculate adaptive correspondence gate based on positional covariance
        pos_variance = trace(cov_pred[:2, :2])
        #  3-sigma radius, clamped between 0.5m and 2.0m
        search_radius: float = float(clip(3.0 * sqrt(pos_variance), 0.5, 2.0))

        for _ in range(max_iter):
            # Transform local scan points into the global map frame
            theta: float = z[2]
            R_curr: ndarray = array([[cos(theta), -sin(theta)], [sin(theta), cos(theta)]])
            global_pts: ndarray = (R_curr @ local_pts.T).T + z[:2]

            # Find the closest Map points (Q) for every Global Scan point (P)
            distances, indices = map_data['kdtree'].query(global_pts)

            # Reject outlier points using the adaptive search radius
            good_matches: ndarray = distances < search_radius
            if sum(good_matches) < 10:
                break  # Lost tracking

            # SVD-Based Alignment (Arun et al., 1987 / Umeyama, 1991)
            P: ndarray = global_pts[good_matches]
            Q: ndarray = map_data['points'][indices[good_matches]]

            # Compute Centroids
            centroid_P: ndarray = mean(P, axis=0)
            centroid_Q: ndarray = mean(Q, axis=0)

            # Center the points around the origin
            P_centered: ndarray = P - centroid_P
            Q_centered: ndarray = Q - centroid_Q

            # Compute Cross-Covariance Matrix (H)
            H: ndarray = P_centered.T @ Q_centered

            # Singular Value Decomposition
            U, _, Vt = svd(H)

            # Optimal Rotation (R_delta)
            R_delta: ndarray = Vt.T @ U.T
            # Ensure it is a pure rotation, not a reflection
            if det(R_delta) < 0:
                Vt[-1, :] *= -1
                R_delta = Vt.T @ U.T

            # Optimal Translation (t_delta)
            t_delta: ndarray = centroid_Q - R_delta @ centroid_P

            # Apply the delta transformation to current pose guess
            z[:2] = R_delta @ z[:2] + t_delta
            z[2] += arctan2(R_delta[1, 0], R_delta[0, 0])

            # Check for convergence based on Mean Squared Error
            mean_error: float = mean(distances[good_matches])
            if abs(prev_error - mean_error) < tolerance:
                break
            prev_error = mean_error

        # Format final pose
        z = wrap_angle(z)

        is_valid: bool = EKFLocalizationNode.inlier_fraction_check(
            current_pose=z, local_pts=local_pts, map_data=map_data
        )

        return z, is_valid, prev_error

    @staticmethod
    def inlier_fraction_check(
        current_pose: ndarray,
        local_pts: ndarray,
        map_data: dict,
        distance_threshold: float = 0.15,
        valid_threshold: float = 0.8,
    ) -> bool:
        # Project scan points to the global map frame using the pose
        final_theta: float = current_pose[2]
        R_final: ndarray = array(
            [[cos(final_theta), -sin(final_theta)], [sin(final_theta), cos(final_theta)]]
        )
        final_global_pts: ndarray = (R_final @ local_pts.T).T + current_pose[:2]

        # Query distances to the nearest map points.
        final_distances, _ = map_data['kdtree'].query(final_global_pts)

        # Calculate the fraction of points within distance_threshold
        inliers = sum(final_distances <= distance_threshold)
        inlier_fraction = float(inliers) / len(final_global_pts)

        # Consider the match "valid" if within valid_threshold
        return inlier_fraction >= valid_threshold

    @staticmethod
    def mapped_process_noise(x: ndarray, M: ndarray, dt: float, q_yaw_scale: float) -> ndarray:
        """R"""
        theta: float = x[2]
        G = array([[cos(theta), 0.0], [sin(theta), 0.0], [0.0, 1.0], [0.0, 0.0]])
        R: ndarray = G @ M @ G.T * dt
        R[3, 3] += q_yaw_scale**2 * dt
        return R

    @staticmethod
    def _predict(
        prev_x: ndarray,
        prev_cov: ndarray,
        v: float,
        omega: float,
        dt: float,
        q_yaw_scale: float,
        alpha_tv_causedby_tv: float,
        alpha_tv_causedby_rv: float,
        alpha_rv_causedby_tv: float,
        alpha_rv_causedby_rv: float,
    ) -> Tuple[ndarray, ndarray]:
        """Prediction step. Apply Kinematic model"""
        # Package the inputs for the motion model
        s: float = prev_x[3]
        u: ndarray = array([v, s * omega])

        # Apply model forward kinematics
        x: ndarray = prev_x.copy()
        x[:3] = UnicycleMotionModel.forward(x=prev_x[:3], u=u, noise=zeros(3), dt=dt)
        x: ndarray = wrap_angle(x)

        # Proportional Control Noise Covariance (M)
        M: ndarray = EKFLocalizationNode.control_noise_covariance(
            u=u,
            alpha_tv_causedby_tv=alpha_tv_causedby_tv,
            alpha_tv_causedby_rv=alpha_tv_causedby_rv,
            alpha_rv_causedby_tv=alpha_rv_causedby_tv,
            alpha_rv_causedby_rv=alpha_rv_causedby_rv,
        )

        # Mapped Process Noise (R)
        R: ndarray = EKFLocalizationNode.mapped_process_noise(
            x=prev_x, M=M, dt=dt, q_yaw_scale=q_yaw_scale
        )

        # Calculate the Jacobian using the previous step.
        A: ndarray = eye(4)
        A[:3, :3] = UnicycleMotionModel.state_jacobian(x=prev_x[:3], u=u, dt=dt)
        A[2, 3] = omega * dt

        # Update the covariance
        cov: ndarray = EKFLocalizationNode.update_covariance(A=A, cov=prev_cov, R=R)

        return x, cov

    @staticmethod
    def _update(
        x_pred: ndarray, cov_pred: ndarray, W_pose: ndarray, z: ndarray, match_error: float
    ) -> Tuple[ndarray, ndarray, bool]:
        error_scale: float = max(1.0, (match_error / 0.05) ** 2)
        W_dynamic: ndarray = W_pose * error_scale

        # Observation matrix (H is identity because we observe the direct [x, y, theta, s] pose from ICP scan matching)
        H: ndarray = eye(3, 4)

        # Innovation covariance (S)
        S: ndarray = H @ cov_pred @ H.T + W_dynamic

        # Innovation (y_tilde) is the difference between observation and prediction
        y_tilde: ndarray = wrap_angle(z - x_pred[:3])

        # Statistical Acceptance Test
        mahalanobis_sq = y_tilde.T @ inv(S) @ y_tilde

        if mahalanobis_sq > chi2.ppf(0.95, 3):
            # The observation is statistically absurd given current belief.
            # Reject the update and keep prior estimate.
            return x_pred, cov_pred, False

        # Kalman Gain (K)
        K: ndarray = cov_pred @ H.T @ inv(S)

        # Update the state estimate and shrink the covariance
        x_update: ndarray = x_pred + K @ y_tilde
        x_update[3] = clip(x_update[3], 0.5, 1.2)
        cov_update: ndarray = (eye(4) - K @ H) @ cov_pred

        # Wrap heading to [-pi, pi] to prevent runaway angle values over time
        x_update = wrap_angle(x_update)

        return x_update, cov_update, True


def main(args=None):
    rclpy.init(args=args)
    node = EKFLocalizationNode()
    try:
        rclpy.spin(node)
    except KeyboardInterrupt:
        pass
    finally:
        node.destroy_node()
        rclpy.shutdown()


if __name__ == '__main__':
    main()
\end{verbatim}

\newpage
\subsection{Particle Filter Node Implementation}\label{subsec:PF implementation}
\begin{verbatim}
#!/usr/bin/env python3

from typing import Optional
from time import perf_counter
from scipy.ndimage import distance_transform_edt
from numpy import (
    ndarray,
    array,
    cov,
    float32 as np_float32,
    ones_like,
    ones,
    searchsorted,
    cumsum,
    average,
    arctan2,
    cos,
    sin,
    copy as np_copy,
    arange,
    exp as np_exp,
    max as np_max,
    zeros,
    vstack,
    where,
    pi,
    mean as np_mean,
)

import rclpy
from rclpy.qos import DurabilityPolicy, QoSProfile, ReliabilityPolicy

# ROS 2 Messages
from nav_msgs.msg import OccupancyGrid, Odometry
from sensor_msgs.msg import LaserScan
from geometry_msgs.msg import PoseArray, Pose

# Import devol localization classes
from devol_localization.base_localization_node import BaseLocalizationNode
from devol_localization.laser_model import LaserModel
from devol_localization.motion_models import UnicycleMotionModel
from devol_localization.utils import euler_to_quaternion


__author__ = 'Jacob Taylor Cassady'
__email__ = 'jcassad1@jh.edu'


class ParticleFilterNode(BaseLocalizationNode):
    def __init__(self):
        super().__init__(
            'particle_filter_node',
            '/devol_drive/pf_pose',
            runtime_topic='/devol_drive/pf_pose/runtime',
        )

        self.declare_parameter('num_particles', '500')

        self.map_dict = None
        self.num_particles: int = int(self.get_parameter('num_particles').value)

        # Initialize particles clustered around starting origin
        self.particles: ndarray = self.rng.multivariate_normal(
            self.x_est, self.covariance, self.num_particles
        )
        self.weights: ndarray = ones(self.num_particles) / self.num_particles

        # Adaptive injection moving averages
        self.slow_avg_lik = 0.0
        self.fast_avg_lik = 0.0

        # Total runtime accumulator
        self.localization_runtime = 0.0

        # Subscribers / Publishers
        map_qos: QoSProfile = QoSProfile(
            depth=1,
            reliability=ReliabilityPolicy.RELIABLE,
            durability=DurabilityPolicy.TRANSIENT_LOCAL,
        )
        self.create_subscription(OccupancyGrid, self.map_topic, self.map_callback, map_qos)
        self.create_subscription(LaserScan, self.scan_topic, self.scan_callback, 1)
        self.create_subscription(Odometry, self.odom_topic, self.odom_callback, 10)

        self.pub_particles = self.create_publisher(PoseArray, self.pose_topic + '_cloud', 10)

        self.get_logger().info('Particle Filter Node Started!')

    def odom_callback(self, msg: Odometry):
        """Base callback to track the latest odometry twist."""
        start_time: float = perf_counter()
        self.v: float = msg.twist.twist.linear.x
        self.omega: float = msg.twist.twist.angular.z
        self.u: ndarray = array([self.v, self.omega])

        dt: Optional[float] = self.get_dt_from_header(msg.header)
        if dt is None:
            return

        self.particles = self._predict(dt=dt)
        self.localization_runtime += perf_counter() - start_time

    def map_callback(self, msg: OccupancyGrid):
        """Converts ROS OccupancyGrid into a distance transform map for the Likelihood field"""
        start_time: float = perf_counter()

        if not self.is_map_new(msg):
            self.localization_runtime += perf_counter() - start_time
            return

        # Convert 1D array to 2D numpy array
        grid: ndarray = array(msg.data, dtype=np_float32).reshape(
            (msg.info.height, msg.info.width)
        )

        # Create binary map where obstacle = 0, free space = 1
        binary_map: ndarray = ones_like(grid)
        binary_map[grid > 50] = 0

        # Calculate Euclidean distance to nearest obstacle (in pixels)
        distance_grid: ndarray = distance_transform_edt(binary_map)
        edt_map: ndarray = distance_grid * msg.info.resolution

        # Pre-compute free space coordinates for the escape hatch
        free_y, free_x = where(binary_map == 1)
        res: float = msg.info.resolution
        ox: float = msg.info.origin.position.x
        oy: float = msg.info.origin.position.y

        # Convert pixel indices to real-world coordinates
        free_world_x: ndarray = free_x * res + ox + (res / 2.0)
        free_world_y: ndarray = free_y * res + oy + (res / 2.0)

        # Store in dict to pass to LaserModel easily
        self.map_dict: dict = {
            'grid': edt_map,
            'info': msg.info,
            'free_x': free_world_x,
            'free_y': free_world_y,
        }
        self.get_logger().info('Map received and Euclidean Distance Transform computed.')
        self.localization_runtime += perf_counter() - start_time

    def scan_callback(self, msg: LaserScan):
        """The main loop of the PF: computes dt, downsamples scan, runs PF, and publishes results."""
        start_time: float = perf_counter()

        # Wait for map and ensure odometry has ran at least once.
        if self.map_dict is None:
            self.localization_runtime += perf_counter() - start_time
            return

        z_dists, z_thetas = self.preprocess_scan(msg=msg)

        if len(z_dists) == 0:
            self.localization_runtime += perf_counter() - start_time
            return  # No valid measurements

        # Run the Particle Filter using Odometry values and the Unicycle model
        self.particles, self.x_est = self._update(z_dists=z_dists, z_thetas=z_thetas)
        self.localization_runtime += perf_counter() - start_time

        # Publish state and particles
        self.publish_state(msg.header.stamp)

    def publish_state(self, stamp):
        """Publishes the estimated pose and the particle cloud for RViz"""
        cov_3x3: ndarray = cov(self.particles, rowvar=False)

        # Publish state
        self.publish_pose_and_tf_runtime(
            stamp=stamp, cov_3x3=cov_3x3, start_time=0.0, end_time=self.localization_runtime
        )

        # Publish Particle Cloud
        cloud_msg: PoseArray = PoseArray()
        cloud_msg.header.stamp = stamp
        cloud_msg.header.frame_id = self.map_frame
        for p in self.particles:
            pose = Pose()
            pose.position.x = p[0]
            pose.position.y = p[1]
            pose.orientation = euler_to_quaternion(p[2])
            cloud_msg.poses.append(pose)

        self.pub_particles.publish(cloud_msg)

        # Reset runtime for the next cycle
        self.localization_runtime = 0.0

    def _predict(self, dt: float) -> ndarray:
        ## Apply the motion model to the samples (Prediction Step)
        # Create noise proportional to the velocity commands for each particle
        sigma_v: float = self.alpha_tv_causedby_tv * abs(self.v) + self.alpha_tv_causedby_rv * abs(
            self.omega
        )
        sigma_w: float = self.alpha_rv_causedby_tv * abs(self.v) + self.alpha_rv_causedby_rv * abs(
            self.omega
        )

        # Add a small floor (0.02) to heading noise to maintain diversity when stopped
        sigma_w = max(0.02, sigma_w)

        # Draw noise applied to the distance and heading moved using sqrt(dt)
        sqrt_dt = dt**0.5

        # Calculate nominal distances and headings moved in this timestep
        dist_moved = self.v * dt
        heading_moved = self.omega * dt

        # Vectorized generation of noisy movements (std_dev = sigma * sqrt(dt))
        noisy_dists = self.rng.normal(dist_moved, sigma_v * sqrt_dt, self.num_particles)
        noisy_headings = self.rng.normal(heading_moved, sigma_w * sqrt_dt, self.num_particles)

        # Convert back to effective velocities since UnicycleMotionModel expects velocity inputs
        noisy_vs = noisy_dists / dt
        noisy_ws = noisy_headings / dt

        # Apply the forward model using the noisy velocities for each particle
        new_particles: ndarray = array(
            [
                UnicycleMotionModel.forward(
                    self.particles[i],
                    array([noisy_vs[i], noisy_ws[i]]),
                    zeros(3),
                    dt,
                )
                for i in range(self.num_particles)
            ]
        )

        # Inject independent X and Y lateral slip noise directly to the array
        new_particles[:, 0] += self.rng.normal(0.0, 0.005, self.num_particles)
        new_particles[:, 1] += self.rng.normal(0.0, 0.005, self.num_particles)

        return new_particles

    def _update(self, z_dists: list, z_thetas: list):
        ## Update
        # Calculate where the laser actually is for each newly predicted particle
        laser_poses: ndarray = np_copy(self.particles)
        laser_offset: ndarray = self.get_laser_offset()
        # Shift X and Y based on the laser's physical mount and the robot's current yaw
        laser_poses[:, 0] += laser_offset[0] * cos(laser_poses[:, 2]) - laser_offset[1] * sin(
            laser_poses[:, 2]
        )
        laser_poses[:, 1] += laser_offset[0] * sin(laser_poses[:, 2]) + laser_offset[1] * cos(
            laser_poses[:, 2]
        )
        # Shift Yaw
        laser_poses[:, 2] += laser_offset[2]

        # Computes log probabilities
        log_weights = array(
            [
                LaserModel.likelihood_field_range_finder(
                    z_dists=z_dists, z_thetas=z_thetas, x=particle, map_dict=self.map_dict, std=0.1
                )
                for particle in laser_poses
            ],
            dtype=float,
        )

        # shift all weights so their highest value is 0.
        max_log_weights: ndarray = np_max(log_weights)
        step_weights: ndarray = np_exp(log_weights - max_log_weights)

        # Calculate raw likelihood average for the adaptive injection logic
        avg_raw_likelihood = np_mean(np_exp(log_weights / len(z_dists)))

        # Initialize running averages on the very first update
        if self.slow_avg_lik == 0.0:
            self.slow_avg_lik = avg_raw_likelihood
            self.fast_avg_lik = avg_raw_likelihood

        # Update exponential moving averages (standard AMCL alpha values)
        self.slow_avg_lik += 0.001 * (avg_raw_likelihood - self.slow_avg_lik)
        self.fast_avg_lik += 0.1 * (avg_raw_likelihood - self.fast_avg_lik)

        # Multiply the new step weights into the running weight accumulator
        self.weights *= step_weights

        # Prevent division by zero
        self.weights += 1e-40
        # normalize (\eta)
        self.weights /= sum(self.weights)

        # Calculate Effective Sample Size (N_eff)
        n_eff = 1.0 / sum(self.weights**2)

        # Estimate the state based on the weighted average of particles
        x_est: ndarray = average(self.particles, axis=0, weights=self.weights)

        # Wrap angle
        x_est[2] = arctan2(
            average(sin(self.particles[:, 2]), weights=self.weights),
            average(cos(self.particles[:, 2]), weights=self.weights),
        )

        # Conditional Resampling: Only resample if N_eff drops below 50%
        if n_eff < (self.num_particles / 2.0):
            # Calculate adaptive escape fraction
            escape_fraction = 0.0
            if self.slow_avg_lik > 0 and self.fast_avg_lik < self.slow_avg_lik:
                # The fast average dropped below the slow average; spike injection rate.
                escape_fraction = 1.0 - (self.fast_avg_lik / self.slow_avg_lik)

            # Cap the maximum injection
            escape_fraction = max(0.0, min(0.05, escape_fraction))

            num_random = int(self.num_particles * escape_fraction)
            num_normal = self.num_particles - num_random
            normal_particles: ndarray = self.low_variance_resample(
                self.particles, self.weights, num_samples=num_normal
            )

            # Draw fresh random particles uniformly from map freespace
            free_x = self.map_dict.get('free_x', array([]))
            free_y = self.map_dict.get('free_y', array([]))

            if num_random > 0:
                if len(free_x) > 0:
                    rand_indices = self.rng.choice(len(free_x), size=num_random, replace=True)
                    random_particles = zeros((num_random, 3))
                    random_particles[:, 0] = free_x[rand_indices]
                    random_particles[:, 1] = free_y[rand_indices]
                    random_particles[:, 2] = self.rng.uniform(-pi, pi, num_random)
                else:
                    # Safe fallback if map processing missed freespace coordinates
                    random_particles = self.rng.multivariate_normal(
                        x_est, cov(self.particles, rowvar=False), num_random
                    )

                # Combine both groups to form the new generation
                new_particles: ndarray = vstack((normal_particles, random_particles))
            else:
                new_particles = normal_particles

            # Reset weights to uniform after a resample event
            self.weights = ones(self.num_particles) / self.num_particles
        else:
            new_particles = self.particles

        return new_particles, x_est

    def low_variance_resample(
        self, particles: ndarray, weights: ndarray, num_samples: Optional[int] = None
    ) -> ndarray:
        N = len(particles) if num_samples is None else num_samples

        # Generate a single random start point between 0 and 1/N
        step = 1.0 / N
        pointers = self.rng.uniform(0, step) + arange(N) * step

        # Find where those evenly spaced pointers land in the cumulative weight array
        cumulative_sum = cumsum(weights)
        indexes = searchsorted(cumulative_sum, pointers)

        return particles[indexes]


def main(args=None):
    rclpy.init(args=args)
    node = ParticleFilterNode()
    try:
        rclpy.spin(node)
    except KeyboardInterrupt:
        pass
    finally:
        node.destroy_node()
        rclpy.shutdown()


if __name__ == '__main__':
    main()
\end{verbatim}
